%% toc: P15 166 166.7
\chapter{Existence, Uniqueness, Multipliers and Reference-Measure Dependence}
\label{sec:p15-w02b-p15-044-maxent-existence}

The general maximum-entropy problem is set up and solved in
\cref{sec:p15-w02b-p15-043-maxent}. There the problem is declared
(\cref{def:p15-w02b-p15-043-problem}), the Gibbs identity and dual bound
are proved (\cref{thm:p15-w02b-p15-043-gibbs}), a feasible
exponential-family member is shown to be the unique maximizer
(\cref{thm:p15-w02b-p15-043-optimality}), and multipliers are shown to
exist and be unique for a regular family with minimal statistics
(\cref{thm:p15-w02b-p15-043-existence}). The same chapter treats the finite
case including boundary values (\cref{thm:p15-w02b-p15-043-finite}), the
ways existence fails (\cref{thm:p15-w02b-p15-043-failures}),
normalization (\cref{sec:p15-w02b-p15-043-multipliers}), and
reparametrization and constant rescaling of the reference measure
(\cref{thm:p15-w02b-p15-043-transport}). The entropy shift under a change
of reference measure is \cref{thm:p15-w02b-p15-042-reference-change}.

This chapter adds three things those results leave implicit. First, a
criterion that decides from the dual function alone whether a finite
supremum is attained. It also shows that an interior minimizer of the dual
function certifies the maximizer without regularity or minimality.
Second, the effect on the \emph{maximizer} of replacing the reference
measure by one with a nonconstant density. Third, the factory example of
H\'ajek and Staffel, computed with and without a constraint on the mean.

\emph{Inputs.} A declared problem in the sense of
\cref{def:p15-w02b-p15-043-problem}, a candidate multiplier or a dual
minimizer, and, for the second part, a positive density relating two
reference measures. \emph{Conclusions.} Whether a maximizer exists and
which one it is; whether a change of reference measure changes the
maximizer, and how its multipliers move when it does not; and, for the
factory, the exact optimization problem behind each of three answers.

\section{Notation}
\label{sec:p15-w02b-p15-044-notation}

The notation is that of
\cref{def:p15-w02b-p15-043-problem,def:p15-w02b-p15-043-partition}: sample
space \(\mathcal X\), reference measure \(\nu\), statistics
\(f=(f_1,\dots,f_m)\), declared values \(c\), feasible set
\(\mathcal F_\nu(c)\), entropy \(H_\nu\), value \(v_\nu(c)\), partition
function \(Z_\nu(\lambda)=\int e^{-\lambda\cdot f}\,d\nu\) with domain
\(\Lambda_\nu\), family members \(Q_\lambda\) with
\(\nu\)-density \(\exp\{-\lambda_0-\lambda\cdot f\}\),
\(\lambda_0=\ln Z_\nu(\lambda)\), and dual function
\(\Psi_c(\lambda)=\ln Z_\nu(\lambda)+\lambda\cdot c\). The sign of the
multipliers follows Jaynes's form, as in that chapter. Logarithms are
natural. Lebesgue measure on an interval is written \(d\ell\).

\section{Attainment and interior dual minimizers}
\label{sec:p15-w02b-p15-044-attainment}

\begin{hbtheorem}{thm:p15-w02b-p15-044-attainment}{proved=annex:p15-w02b-p15-044-attainment}{Attainment criterion and certification by an interior dual minimizer}
Consider a declared problem with \(\Lambda_\nu\neq\varnothing\).
\begin{enumerate}
\item Suppose \(v_\nu(c)=\Psi_c(\lambda^\circ)\) for some
\(\lambda^\circ\in\Lambda_\nu\). Then \(\lambda^\circ\) minimizes
\(\Psi_c\) on \(\Lambda_\nu\), and a maximizer exists if and only if
\(Q_{\lambda^\circ}\in\mathcal F_\nu(c)\); in that case the maximizer is
\(Q_{\lambda^\circ}\).
\item If \(\lambda\in\Lambda_\nu\) and \(Q_\lambda\in\mathcal F_\nu(c)\),
then \(v_\nu(c)=\Psi_c(\lambda)=\min_{\Lambda_\nu}\Psi_c\).
\item If \(\lambda^*\) is an interior point of \(\Lambda_\nu\) and
minimizes \(\Psi_c\) on \(\Lambda_\nu\), then
\(Q_{\lambda^*}\in\mathcal F_\nu(c)\), \(Q_{\lambda^*}\) is the unique
maximizer, and \(v_\nu(c)=\Psi_c(\lambda^*)\). Neither regularity of the
family nor minimality of the statistics is needed.
\end{enumerate}
\end{hbtheorem}

Part~1 turns a duality statement into an existence test. Once the value is
known to equal the dual function at some multiplier, only one distribution
can attain it, and checking its constraints settles existence. Item~3 of
\cref{thm:p15-w02b-p15-043-failures} is an instance in which the check
fails: there \(v_\nu(c)=\Psi_c(0)\), and \(Q_0\) has the wrong mean, so no
maximizer exists. Part~3 separates two roles that regularity and
minimality play in \cref{thm:p15-w02b-p15-043-existence}. They guarantee
that a solution of the multiplier equations exists and is unique. They are
not needed to show that a solution, once found at an interior point, gives
the maximizer. With redundant statistics, such as \(f_2=2f_1\), the
interior minimizers of \(\Psi_c\) form a line segment; part~3 applies to
each point of it, and all give the same \(Q_\lambda\), as item~4 of
\cref{thm:p15-w02b-p15-043-partition} also shows.

A numerical routine that minimizes \(\Psi_c\) and returns a multiplier on
the edge of \(\Lambda_\nu\), or a multiplier sequence of growing size,
has not certified anything. Parts~1 and~2 say exactly what must still be
checked: that the returned \(Q_\lambda\) satisfies the declared
constraints.

\section{Changing the reference measure by a density}
\label{sec:p15-w02b-p15-044-reference}

\Cref{thm:p15-w02b-p15-043-transport} shows that relabelling the
outcomes changes nothing when the reference measure is carried along, and
that multiplying it by a constant does not change the maximizer. The case
that matters for representation dependence is a reference measure
redeclared with a nonconstant density.

\begin{hbtheorem}{thm:p15-w02b-p15-044-reference}{proved=annex:p15-w02b-p15-044-reference}{Reference change: a fixed-multiplier statistic and an affine test}
Let \(\nu'=\rho\,\nu\), where \(\rho:\mathcal X\to(0,\infty)\) is
measurable and \(\nu'\) is \(\sigma\)-finite. The two problems with the
same statistics \(f\) and values \(c\) have the same feasible set.
Write \(Z'\), \(\Lambda'\), \(Q'_\lambda\) and \(\Psi'_c\) for the objects
of the problem with reference \(\nu'\).
\begin{enumerate}
\item \(Z'(\lambda)=\int e^{-\lambda\cdot f+\ln\rho}\,d\nu\). Hence
\(\lambda\in\Lambda'\) exactly when \((\lambda,-1)\) lies in the
partition-function domain of the problem with reference \(\nu\) and
statistics \((f,\ln\rho)\), and \(Q'_\lambda\) is that problem's family
member with multipliers \((\lambda,-1)\).
\item If \(\ln\rho=b+a\cdot f\) \(\nu\)-almost everywhere for some
\(b\in\mathbb R\) and \(a\in\mathbb R^m\), then
\(\Lambda'=\Lambda_\nu+a\), \(Q'_\lambda=Q_{\lambda-a}\) and
\(\Psi'_c(\lambda)=\Psi_c(\lambda-a)+b+a\cdot c\). For every feasible
\(P\), \(H_\nu(P)\) is defined exactly when \(H_{\nu'}(P)\) is, and then
\(H_{\nu'}(P)=H_\nu(P)+b+a\cdot c\). The two problems have the same
maximizers, and their multipliers differ by \(a\).
\item If \(\ln\rho\) is not \(\nu\)-almost everywhere equal to an affine
function \(b+a\cdot f\), then no \(Q'_\lambda\) equals any \(Q_\mu\).
Consequently, if each problem has a maximizer belonging to its own family,
the two maximizers differ.
\end{enumerate}
\end{hbtheorem}

Part~1 is the precise sense in which a new reference measure is a hidden
constraint. Declaring \(\nu'\) instead of \(\nu\) has the same effect on
the family as adding the statistic \(\ln\rho\) and fixing its multiplier in
advance, without declaring any value for its expectation. Parts~2 and~3
give the test: the answer is unchanged when \(\ln\rho\) is already an
affine function of the declared statistics, and it changes, whenever both
problems are solved inside their families, when it is not.

\section{The factory example with a declared mean}
\label{sec:p15-w02b-p15-044-cubes}

\begin{hbtheorem}{thm:p15-w02b-p15-044-cubes}{proved=annex:p15-w02b-p15-044-cubes}{Cube side lengths under three references}
Let \(\mathcal X=(0,1]\) with Borel sets, write \(\ell\) for the side
length, and for \(s=0,1,2\) let \(\nu_s(d\ell)=(s+1)\ell^s\,d\ell\).
\begin{enumerate}
\item \(\nu_s\) is the image of Lebesgue measure on \((0,1]\) under
\(u\mapsto u^{1/(s+1)}\). Thus \(\nu_0\), \(\nu_1\) and \(\nu_2\) are
Lebesgue measure in side length, face area and volume, each written in the
side-length coordinate.
\item With no constraint (\(m=0\)), the unique maximizer for \(\nu_s\) is
\(\nu_s\) itself, which gives the event \(\ell\leq1/2\) probability
\(2^{-(s+1)}\): \(1/2\), \(1/4\) and \(1/8\).
\item With \(m=1\), \(f(\ell)=\ell\) and \(c=1/2\), for each \(s\) there is
exactly one \(\lambda_s\in\mathbb R\) with
\[
   \frac{\int_0^1\ell^{s+1}e^{-\lambda_s\ell}\,d\ell}
        {\int_0^1\ell^{s}e^{-\lambda_s\ell}\,d\ell}=\frac12 ,
\]
and the unique maximizer has Lebesgue density proportional to
\(\ell^se^{-\lambda_s\ell}\). Here \(\lambda_0=0\) and
\(0<\lambda_1,\lambda_2<10\).
\item The three maximizers of part~3 are pairwise distinct. The \(\nu_2\)
maximizer is the member of the \(\nu_0\) family for the statistics
\((\ell,\ln\ell)\) with multipliers \((\lambda_2,-2)\).
\end{enumerate}
\end{hbtheorem}

\begin{hbexample}{def:p15-w02b-p15-044-cube-example}{One declared mean, three answers}
H\'ajek and Staffel use a factory producing cubes with sides of at most one
foot. They show that indifference gives conflicting answers when side
length, face area and volume compete as the quantity over which credence
is spread uniformly \cite[p.~244]{HajekStaffel2021Dynamics}. Part~2 of
\cref{thm:p15-w02b-p15-044-cubes} places their answers inside the declared
problem: \(1/2\), \(1/4\) and \(1/8\) are the maximizers for three
reference measures with no constraint. The source gives \(1/2\) and
\(1/8\) and says only that face area yields yet another answer; the value
\(1/4\) is the computation of part~2, as in \cref{sec:p15-rat41-synchronic}.

Now add the declared premise that the mean side length is \(1/2\) foot.
Solving the equation of part~3 numerically gives
\[
  \lambda_1\approx2.688,\qquad \lambda_2\approx5.071,
\]
and the probabilities of \(\ell\leq1/2\) become
\[
  0.5000\ (\nu_0),\qquad 0.5188\ (\nu_1),\qquad 0.5280\ (\nu_2).
\]
Each figure is the exact solution of its declared problem. The declared
mean narrows the spread of answers from \([1/8,1/2]\) to
\([0.500,0.528]\), but it does not remove it, and by part~4 it cannot. By
\cref{thm:p15-w02b-p15-044-reference}, the volume reference acts on the
side-length problem as the extra statistic \(\ln\ell\) with its multiplier
fixed at \(-2\), because \(\ln(3\ell^2)=\ln3+2\ln\ell\) is not affine in
\(\ell\). Nothing in the optimization indicates which reference measure the
available evidence supports. That is a premise, argued for or against on
other grounds, and the declared mean is a premise of the same kind.
\end{hbexample}

\section{Procedure and limits}
\label{sec:p15-w02b-p15-044-procedure}

\begin{hbmethod}{def:p15-w02b-p15-044-procedure}{Certifying a candidate and testing a change of reference}
This supplements \cref{def:p15-w02b-p15-043-method}.
\begin{enumerate}
\item For a candidate multiplier \(\lambda\), check
\(\lambda\in\Lambda_\nu\), the integrability of each \(f_k\) under
\(Q_\lambda\), and \(E_{Q_\lambda}[f_k]=c_k\). If all hold, \(Q_\lambda\)
is the maximizer by part~2 of \cref{thm:p15-w02b-p15-044-attainment}.
\item If the dual minimum is found at an interior point of
\(\Lambda_\nu\), the checks of step~1 hold automatically, by part~3 of
that theorem.
\item If the value equals \(\Psi_c(\lambda^\circ)\) with \(\lambda^\circ\)
on the edge of \(\Lambda_\nu\), decide existence by checking whether
\(Q_{\lambda^\circ}\) is feasible (part~1); if it is not, report that the
value is not attained.
\item For a proposed new reference \(\nu'=\rho\nu\), compute \(\ln\rho\).
If it is affine in the declared statistics, the maximizer is unchanged and
the multipliers shift by \(a\). Otherwise solve both problems and report
both answers together with their reference measures
(\cref{thm:p15-w02b-p15-044-reference}).
\end{enumerate}
\end{hbmethod}

\begin{hbwarning}{def:p15-w02b-p15-044-limits}{What the tests do not settle}
\begin{enumerate}
\item The test of \cref{thm:p15-w02b-p15-044-reference} compares
maximizers only when both lie in their exponential families. At a boundary
value in the finite case the maximizer lies outside its family
(\cref{thm:p15-w02b-p15-043-finite}), and a reference change with
nonaffine \(\ln\rho\) can then leave it unchanged. On \(\{1,2,3\}\) with
\(f(x)=x\) and \(c=1\), the only feasible distribution is the point mass
at \(1\), so it is the maximizer for every choice of positive weights.
\item The optimization results are conditional on the declared \(\nu\),
\(f\) and \(c\). Whether the maximizer should be adopted as a credence is
a question about the inferential argument of
\cref{sec:p15-w02b-p15-043-jaynes}. H\'ajek and Staffel record the
representation objection as disputed, citing Norton, White and Novack
for defences of indifference \cite[p.~244]{HajekStaffel2021Dynamics}.
\item The numerical values in \cref{def:p15-w02b-p15-044-cube-example}
are rounded solutions of the stated equations. The theorem proves the
existence, uniqueness and location in \((0,10)\) of the multipliers, not
their decimal expansions.
\end{enumerate}
\end{hbwarning}

\section{Relationships and sources}
\label{sec:p15-w02b-p15-044-relations}

\Cref{thm:p15-w02b-p15-044-attainment} uses only the Gibbs identity and
the partition-function lemma of \cref{sec:p15-w02b-p15-043-maxent}.
\Cref{thm:p15-w02b-p15-044-reference} extends the constant rescaling of
\cref{thm:p15-w02b-p15-043-transport} to arbitrary positive densities,
using the entropy shift of \cref{thm:p15-w02b-p15-042-reference-change}.
\Cref{def:p15-w02b-p15-043-reference-example} compares the references
\(dx\) and \(x\,dx\) on \((0,\infty)\) under a declared mean; the factory
example above does the same on a bounded interval for the three references
named in the source. The representation passage of
\cref{sec:p15-rat41-synchronic} and the reparameterization example of
\cref{thm:p15-bayes-transformation} concern the same phenomenon for priors
chosen by symmetry. With no constraint and unnormalized weights
\((0.3,0.4)\), \cref{thm:p15-w02b-p15-043-no-constraint} gives
\((3/7,4/7)\), the minimizer of \(\sum_iq_i\log_2(q_i/c_i)\) in
\cref{thm:p15-rat41-graded}.

Sources. The factory example and the maximum-entropy generalization of
indifference are from H\'ajek and Staffel
\cite[p.~244]{HajekStaffel2021Dynamics}. The attainment criterion, the
reference-change results and the computations of the factory example are
proved in the annex.
